\chapter{Deployment, Observability und Sicherheit}
\label{chap:deployment_observability}

% ============================================================
% AUTOR-NOTIZ STATT LERNZIELE
% ============================================================
\autornotiz{
2023 habe ich mein erstes LLM-System in Produktion deployed. FastAPI, Docker, OpenAI API.
Kein Caching, keine Retry-Logik, kein Circuit Breaker. Erster Traffic-Spike: API-Latenz 15s,
Error Rate 40\%, Kosten explodiert. Der Fix: Redis-Cache (30% Cost Reduction), Tenacity Retry
mit Exponential Backoff, Circuit Breaker Pattern, Prometheus Metriken + Grafana Dashboard.
Die Lektion: LLM-Deployment ist nicht wie klassisches Web-Deployment. External Dependencies,
probabilistisches Verhalten und Kosten pro Request erfordern neue Patterns.
}

% ============================================================
% MOTIVATION
% ============================================================
\section{Motivation}

Eine LLM-Anwendung in Produktion zu bringen unterscheidet sich grundlegend vom klassischen Software-Deployment. In einer traditionellen Web-App ist das Verhalten deterministisch: Gleicher Input, gleicher Output. Bei LLMs nicht -- das hat weitreichende Konsequenzen.

Drei Herausforderungen machen LLM-Deployment einzigartig:

\begin{enumerate}[leftmargin=*]
    \item \textbf{External Dependencies}: Deine App hängt von der API-Verfügbarkeit ab (OpenAI, Anthropic). Fällt die API aus, fällt deine App mit aus. Keine direkte Kontrolle über Latenz oder Verfügbarkeit.
    \item \textbf{Kosten-Kontrolle}: Jeder Request kostet Geld. Ohne Monitoring explodieren die Kosten bei Popularität. Ein Traffic-Spike kann innerhalb von Stunden tausende Euro kosten.
    \item \textbf{Probabilistisches Verhalten}: Dein Test kann heute bestehen und morgen trotzdem falsche Antworten liefern. Deterministische Unit-Tests allein reichen nicht -- kontinuierliche Evaluation in Produktion ist nötig.
\end{enumerate}

\praxishinweis{External Dependencies sind kein reines Availability-Problem. Ein API-Provider kann verfügbar sein, aber plötzlich 10x höhere Latenz liefern. Baue Timeout-Monitoring und Latenz-Alarme ein -- nicht nur Health-Checks.}

% ============================================================
% GRUNDLAGEN
% ============================================================
\section{Grundlagen -- Deployment-Arten für LLMs}

Vier grundlegende Deployment-Modelle für LLM-Anwendungen:

\begin{itemize}

    \item \textbf{API-basiert (Managed)} Externe API (OpenAI, Anthropic, Google). Kein Modell-Hosting. Pay-per-Token. Einfachster Einstieg, aber Vendor-Lock-in und Latenz-abhängig.
    \item \textbf{Self-hosted Open Source} Open-Source-Modell (Llama, Mistral, Qwen) auf eigener Infrastruktur. Volle Kontrolle, keine Daten verlassen den Server. Hohe Infrastruktur-Kosten (GPUs).
    \item \textbf{Managed Inference (Cloud)} Inference-Anbieter (Together AI, Fireworks, Groq) bietet Open-Source-Modelle als API. Kompromiss aus Kontrolle und Einfachheit.
    \item \textbf{Edge Deployment} Modell läuft auf Endgerät (Smartphone, Laptop). Minimale Latenz, keine Internetverbindung nötig. Stark limitierte Modellgrösse.
\end{itemize}

\subsection{Entscheidungsmatrix}

\begin{table}[H]
\centering
\begin{tabularx}{\linewidth}{lXXXX}
\toprule
\textbf{Kriterium} & \textbf{API} & \textbf{Self-hosted} & \textbf{Cloud Inference} & \textbf{Edge} \\
\midrule
Setup-Aufwand & Sehr gering & Sehr hoch & Gering & Mittel \\
Latenz & 200--2000ms & 50--500ms & 100--1000ms & 10--100ms \\
Kosten pro Query & Mittel--Hoch & Niedrig (fix) & Niedrig--Mittel & Sehr gering \\
Kontrolle & Gering & Voll & Mittel & Voll \\
Datenhoheit & Extern & Intern & Extern & Lokal \\
Skalierung & Automatisch & Manuell & Automatisch & Gerätebasiert \\
\bottomrule
\end{tabularx}
\caption{Vergleich der Deployment-Modelle für LLMs}
\end{table}

% ============================================================
% ARCHITEKTUR
% ============================================================
\subsection{Build-vs-Buy: Wann was wo deployen?}

Die eigentliche Frage ist nicht "welches Modell" sondern "welches Deployment-Modell". Die folgende Entscheidungsmatrix hilft dir, die richtige Wahl basierend auf drei Faktoren zu treffen: Datenhoheit, Scale und Latenz.

\begin{table}[H]
\centering
\small
\begin{tabularx}{\linewidth}{l X X X}
\toprule
\textbf{Szenario} & \textbf{Empfohlen} & \textbf{Warum} & \textbf{Risiko} \\
\midrule
Startups \& MVP & Cloud API (OpenAI, Anthropic) & Kein Infrastruktur-Overhead, sofort einsatzbereit & Vendor-Lock-in, Kostenexplosion bei Scale \\
Mid-Size mit PII-Daten & Self-hosted Llama/Mistral & Daten verlassen niemals den Server & GPU-Kosten,运维-Overhead \\
Geringe Latenz nötig & Edge oder Cloud Inference & Sub-100ms für Edge, \textless{} 200ms für Cloud Inference & Limitierte Modellgrösse (Edge) \\
\textgreater{} 1M Requests/Tag & Self-hosted oder Hybrid & Kosteneffizienz bei Scale & Hohe Initialinvestition \\
Regulierte Industrie & Self-hosted + PII-Filtering & Compliance-Nachweis, vollständige Audit-Trails & Komplexität \\
Hybrid-Architektur & Mix: Cloud + Self-hosted & Kritische Calls self-hosted, Standard-Requests cloud & Architektur-Komplexität \\
\bottomrule
\end{tabularx}
\caption{Build-vs-Buy Decision Tree -- Szenariobasierte Empfehlungen}
\end{table}

\realitycheck{Die beste Entscheidung ist selten "entweder-oder". Die meisten erfolgreichen Systeme verwenden ein Hybrid-Modell: OpenAI für den Standard-Fall, self-hosted Llama für PII-sensitive Calls. Du beginnst mit "Buy", gehst zu "Build" wenn die Kosten es erfordern.}

\praxishinweis{Fang klein an. Deploy mit einer Cloud API. Wenn die Kosten explodieren, wechsle zu einem Hybrid-Modell. Wenn das nicht mehr reicht, gehe zu Self-hosted. Der perfekte Zeitpunkt für "Build" kommt, wenn du die Daten und das Know-how hast.}

\subsection{Architektur einer LLM-Produktionsanwendung}

Eine produktionsreife LLM-Anwendung besteht aus mehreren Schichten:

\begin{verbatim}
+------------------+     +------------------+     +------------------+
|    Frontend      |---->|   API Gateway    |---->|  LLM-Service     |
| (Streamlit,      |     | (Rate Limit,     |     | (FastAPI,        |
|  React, Mobile)  |     |  Auth, Routing)  |     |  Fallback)       |
+------------------+     +------------------+     +--------+---------+
                                                             |
                                               +-------------+-------------+
                                               |             |             |
                                         +-----v------+ +--v--------+ +-v--------+
                                         |  Cache     | |  Vector   | | Tracing  |
                                         |  (Redis)   | |  DB       | |(Langfuse)|
                                         +------------+ +-----------+ +----------+
\end{verbatim}

Jede Schicht hat spezifische Aufgaben:

\begin{enumerate}[leftmargin=*]
    \item \textbf{Frontend}: UI für User-Interaktion, sendet Requests an API Gateway
    \item \textbf{API Gateway}: Authentifizierung, Rate Limiting, Request-Validierung, Routing
    \item \textbf{LLM-Service}: Kernlogik, Prompt-Konstruktion, LLM-Aufruf, Response-Parsing
    \item \textbf{Cache}: Zwischenspeichert häufige Queries, reduziert API-Kosten und Latenz
    \item \textbf{Vector DB}: Speichert und retrievet Embeddings für RAG
    \item \textbf{Tracing}: Protokolliert jeden Request für Debugging und Kosten-Tracking
\end{enumerate}

\praxishinweis{API Gateway und LLM-Service sollten getrennte Deployments sein. Ein Rate-Limit-Treffer soll nicht den gesamten Service blockieren, sondern nur den entsprechenden Client. Trennung erlaubt unabhängiges Skalieren.}

% ============================================================
% THEORIE
% ============================================================
\section{Theorie -- Latenz, Throughput und Kosten}

\subsection{Latenz-Komponenten}

Die Gesamtlatenz eines LLM-Requests setzt sich aus mehreren Komponenten zusammen:

\begin{equation}
T_{total} = T_{network} + T_{queue} + T_{prompt} + T_{generation} + T_{post}
\end{equation}

\begin{itemize}

    \item \textbf{$T_{network}$} Netzwerk-Latenz zwischen Client und API. Bei Cloud-APIs typisch 50--200ms. Standort-Wahl (us-east vs. eu-west) macht bis zu 100ms Unterschied.
    \item \textbf{$T_{queue}$} Wartezeit auf dem Server. Bei hoher Auslastung können Requests Minuten in der Queue hängen. Rate Limiting verstärkt diesen Effekt.
    \item \textbf{$T_{prompt}$} Zeit für die Verarbeitung des Inputs (Pre-fill). Skaliert linear mit der Prompt-Länge. Ein 10K-Token-Prompt dauert ca. 5--15x länger als ein 1K-Token-Prompt.
    \item \textbf{$T_{generation}$} Zeit für die Generierung des Outputs. Skaliert linear mit Output-Tokens und quadratisch mit der Kontextlänge (Attention). Streaming reduziert die Wahrnehmung, nicht die Dauer.
    \item \textbf{$T_{post}$} Post-Processing: Output-Parsing, Safety-Checks, Logging. Typisch < 10ms.
\end{itemize}

\subsection{Das Cost-Performance-Triangle}

\begin{table}[H]
\centering
\begin{tabularx}{\linewidth}{XXXX}
\toprule
\textbf{Strategie} & \textbf{Qualität} & \textbf{Latenz} & \textbf{Kosten} \\
\midrule
Kleineres Modell (4o-mini statt 4o) & -- & ++ & ++ \\
Längerer Prompt (mehr Kontext) & + & -- & -- \\
Caching aktivieren & = & ++ & ++ \\
Streaming aktivieren & = & + (gefühlt) & = \\
Batching von Requests & = & = & ++ \\
Speculative Decoding & = & ++ & = \\
\bottomrule
\end{tabularx}
\caption{Trade-offs zwischen Qualität, Latenz und Kosten}
\end{table}

% ============================================================
% PRAXIS
% ============================================================
\section{FastAPI + Docker -- der Standard-Stack}

Der einfachste Weg, eine LLM-Anwendung als Microservice zu deployen:

\begin{lstlisting}[language=Python, caption=FastAPI-Microservice mit Caching und Rate Limiting]
from fastapi import FastAPI, HTTPException, Request
from fastapi.middleware.cors import CORSMiddleware
from slowapi import Limiter, _rate_limit_exceeded_handler
from slowapi.util import get_remote_address
from slowapi.errors import RateLimitExceeded
import openai
import hashlib
import time
from typing import Optional

app = FastAPI(title="AI-Assistant API", version=1.0.0)
limiter = Limiter(key_func=get_remote_address)
app.state.limiter = limiter
app.add_exception_handler(RateLimitExceeded, _rate_limit_exceeded_handler)

client = openai.OpenAI(api_key=os.environ["OPENAI_API_KEY"])

cache = {}

@app.post("/chat")
@limiter.limit("10/minute")
async def chat(request: Request, user_input: str):
    """Chat-Endpoint mit Caching und Kosten-Tracking."""

    cache_key = hashlib.md5(user_input.encode()).hexdigest()
    if cache_key in cache and time.time() - cache[cache_key]["ts"] < 300:
        return {"response": cache[cache_key]["response"], "cached": True}

    try:
        start = time.time()
        response = client.chat.completions.create(
            model="gpt-4o-mini",
            messages=[
                {"role": "system", "content": get_system_prompt()},
                {"role": "user", "content": user_input},
            ],
            temperature=0.1,
            max_tokens=1024,
        )
        latency_ms = (time.time() - start) * 1000

        result = response.choices[0].message.content

        cache[cache_key] = {
            "response": result,
            "ts": time.time(),
        }

        print(f"[METRIC] latency={latency_ms:.0f}ms tokens={response.usage.total_tokens}")

        return {"response": result, "cached": False, "latency_ms": latency_ms}

    except openai.RateLimitError:
        raise HTTPException(status_code=429, detail="Rate Limit erreicht.")
    except openai.APIConnectionError:
        raise HTTPException(status_code=503, detail="KI-Dienst nicht verfuegbar.")

@app.get("/health")
async def health():
    """Health-Check fuer Load Balancer."""
    return {"status": "healthy", "model": "gpt-4o-mini"}
\end{lstlisting}

\subsection{Docker-Konfiguration}

\begin{lstlisting}[language=bash, caption=Multi-Stage Dockerfile für LLM-Service]
FROM python:3.12-slim AS builder

WORKDIR /app
COPY requirements.txt .
RUN pip install --no-cache-dir -r requirements.txt

FROM python:3.12-slim AS runtime

WORKDIR /app
COPY --from=builder /usr/local/lib/python3.12/site-packages /usr/local/lib/python3.12/site-packages
COPY . .

EXPOSE 8000
CMD ["uvicorn", "main:app", "--host", "0.0.0.0", "--port", "8000"]
\end{lstlisting}

% ============================================================
% LATENZ
% ============================================================
\section{Latenz-Optimierung in der Produktion}

Latenz ist die wichtigste UX-Metrik bei LLM-Anwendungen:

\begin{table}[H]
\centering
\begin{tabularx}{\linewidth}{lXr}
\toprule
\textbf{Technik} & \textbf{Wirkung} & \textbf{Aufwand} \\
\midrule
Streaming-Response & Time-to-first-token < 200ms & Niedrig \\
Caching (Redis) & 30--50\% Reduktion API-Calls & Mittel \\
Modell-Tiering & 4o-mini statt 4o: 2x schneller, 5x günstiger & Niedrig \\
Speculative Decoding & Kleines Modell errät Antwort, grosses validiert & Hoch \\
Batching / Parallel Calls & Mehrere Anfragen gleichzeitig & Mittel \\
Context-Kompaktion & Kürzere Prompts = schnellere Response & Niedrig \\
\bottomrule
\end{tabularx}
\caption{Latenz-Optimierungstechniken}
\end{table}

\subsection{Streaming implementieren}

\begin{lstlisting}[language=Python, caption=Streaming-Response für Echtzeit-Feedback]
from fastapi.responses import StreamingResponse

async def stream_chat(user_input: str):
    """Streaming-Endpoint mit SSE (Server-Sent Events)."""
    stream = client.chat.completions.create(
        model="gpt-4o-mini",
        messages=[{"role": "user", "content": user_input}],
        stream=True,
    )

    async def generate():
        for chunk in stream:
            if chunk.choices[0].delta.content:
                yield f"data: {chunk.choices[0].delta.content}\n\n"
        yield "data: [DONE]\n\n"

    return StreamingResponse(generate(), media_type="text/event-stream")

@app.post("/chat/stream")
async def chat_stream(request: Request, user_input: str):
    return await stream_chat(user_input)
\end{lstlisting}

\praxishinweis{Streaming verbessert die subjektive Latenz massiv, reduziert aber nicht die Gesamtverarbeitungszeit. Der erste Token kommt schneller, die komplette Antwort braucht gleich lang. Für UX zählt Time-to-First-Token mehr als Total-Time.}

% ============================================================
% FEHLER-TOLERANZ
% ============================================================
\section{Fehler-Toleranz und Fallbacks}

LLM-APIs fallen aus. Rate Limits werden erreicht:

\begin{lstlisting}[language=Python, caption=Fehlertolerante LLM-Aufrufe mit Retry und Fallback]
import asyncio
from tenacity import retry, stop_after_attempt,
    wait_exponential, retry_if_exception_type

@retry(
    stop=stop_after_attempt(3),
    wait=wait_exponential(multiplier=1, min=2, max=10),
    retry=retry_if_exception_type(
        (openai.RateLimitError, openai.APIConnectionError)),
    reraise=True,
)
async def robust_chat(messages: list[dict]) -> str:
    """LLM-Aufruf mit automatischem Retry und Fallback."""

    try:
        response = await asyncio.wait_for(
            client.chat.completions.create(
                model="gpt-4o-mini",
                messages=messages,
                temperature=0.1,
                max_tokens=1024,
            ),
            timeout=30.0
        )
        return response.choices[0].message.content

    except openai.RateLimitError:
        # Fallback 1: Modell wechseln
        response = await client.chat.completions.create(
            model="gpt-4o",
            messages=messages,
        )
        return response.choices[0].message.content

    except openai.APIConnectionError:
        # Fallback 2: Kleineres Modell
        response = await client.chat.completions.create(
            model="gpt-4o-mini",
            messages=messages,
            temperature=0.3,
        )
        return response.choices[0].message.content
\end{lstlisting}

\subsection{Circuit-Breaker-Pattern}

Bei wiederholten API-Fehlern verhindert ein Circuit-Breaker weitere Versuche und gibt eine Fallback-Antwort:

\begin{lstlisting}[language=Python, caption=Circuit-Breaker für LLM-API-Ausfälle]
import time
from enum import Enum

class CircuitState(Enum):
    CLOSED = "closed"       # Normalbetrieb
    OPEN = "open"           # API down, keine Versuche
    HALF_OPEN = "half_open" # Test-Request erlaubt

class LLMCircuitBreaker:
    def __init__(self, failure_threshold: int = 5,
                 recovery_timeout: float = 60.0):
        self.state = CircuitState.CLOSED
        self.failure_count = 0
        self.failure_threshold = failure_threshold
        self.recovery_timeout = recovery_timeout
        self.last_failure_time = 0.0

    def call(self, llm_func, *args, **kwargs):
        if self.state == CircuitState.OPEN:
            if time.time() - self.last_failure_time > self.recovery_timeout:
                self.state = CircuitState.HALF_OPEN
            else:
                return self.fallback_response()

        try:
            result = llm_func(*args, **kwargs)
            if self.state == CircuitState.HALF_OPEN:
                self.state = CircuitState.CLOSED
                self.failure_count = 0
            return result
        except Exception as e:
            self.failure_count += 1
            self.last_failure_time = time.time()
            if self.failure_count >= self.failure_threshold:
                self.state = CircuitState.OPEN
            return self.fallback_response()

    def fallback_response(self):
        return {"response": "Service voruebergehend nicht "
                            "verfuegbar. Bitte spaeter erneut versuchen.",
                "cached": False, "fallback": True}

breaker = LLMCircuitBreaker(failure_threshold=3, recovery_timeout=30)
result = breaker.call(robust_chat, [{"role": "user", "content": "Hallo"}])
\end{lstlisting}

% ============================================================
% OBSERVABILITY
% ============================================================
\section{Observability -- Produktion überwachen}

Observability ruht auf drei Säulen:

\begin{itemize}

    \item \textbf{Metriken} Numerische Messwerte über Zeit: Requests/Minute, P95-Latenz, Cost/Query, Error Rate. Instrumentiere jeden Request mit Prometheus-kompatiblen Metriken.
    \item \textbf{Logging} Strukturierte Ereignis-Aufzeichnung: Jeder LLM-Call mit Prompt, Response, Latenz, Kosten. Verwende strukturiertes JSON-Logging für maschinenlesbare Logs.
    \item \textbf{Tracing} Request-Kontext über verteilte Systeme hinweg: Ein User-Request durchläuft API-Gateway, LLM-Service, Vector-DB, Cache. Tracing verbindet diese Spans zu einem Gesamtbild.
\end{itemize}

\begin{lstlisting}[language=Python, caption=Strukturiertes Logging mit Metriken für Prometheus]
import structlog
import time
from prometheus_client import Counter, Histogram, Gauge

# Prometheus-Metriken
LLM_REQUESTS = Counter(
    "llm_requests_total", "Total LLM requests",
    ["model", "status"]
)
LLM_LATENCY = Histogram(
    "llm_latency_seconds", "LLM request latency",
    ["model"], buckets=[0.1, 0.5, 1.0, 2.0, 5.0, 10.0]
)
LLM_COST = Counter(
    "llm_cost_total", "Total LLM cost in EUR",
    ["model"]
)

logger = structlog.get_logger()

def monitored_llm_call(model: str, messages: list) -> str:
    """LLM-Aufruf mit Metriken und Logging."""
    start = time.time()

    try:
        response = client.chat.completions.create(
            model=model, messages=messages
        )
        latency = time.time() - start
        tokens = response.usage.total_tokens
        cost = calculate_cost(model, tokens)

        LLM_REQUESTS.labels(model=model, status="success").inc()
        LLM_LATENCY.labels(model=model).observe(latency)
        LLM_COST.labels(model=model).inc(cost)

        logger.info("llm_call_completed",
            model=model,
            latency_s=round(latency, 3),
            tokens=tokens,
            cost_eur=round(cost, 5),
            prompt_tokens=response.usage.prompt_tokens,
            completion_tokens=response.usage.completion_tokens,
        )

        return response.choices[0].message.content

    except Exception as e:
        latency = time.time() - start
        LLM_REQUESTS.labels(model=model, status="error").inc()
        logger.error("llm_call_failed",
            model=model,
            latency_s=round(latency, 3),
            error=str(e),
        )
        raise
\end{lstlisting}

% ============================================================
% BEST PRACTICES
% ============================================================
\section{Best Practices für LLM-Deployment}

\begin{enumerate}[leftmargin=*]
    \item \textbf{Immer Timeout setzen}: LLM-APIs können minutenlang hängen. Ohne Timeout blockiert dein ganzer Stack. Setze 10--30s Timeout für jeden Request.
    \item \textbf{Caching von Anfang an}: Bei 60\% der Support-Fragen sind die Antworten ähnlich. Ein Redis-Cache spart 30--50\% API-Kosten.
    \item \textbf{Multi-Provider-Strategie}: Hänge dich nicht an einen einzigen API-Anbieter. Baue ein Abstraktions-Layer, das zwischen OpenAI, Anthropic und Open-Source-Modellen wechseln kann.
    \item \textbf{Health-Check implementieren}: Load Balancer brauchen einen Health-Endpoint. Der sollte auch die LLM-API-Verfügbarkeit prüfen.
    \item \textbf{Rate Limiting auf mehreren Ebenen}: API-Gateway-Ebene (pro IP) und LLM-Service-Ebene (pro API-Key). Verhindert Kosten-Explosion bei Traffic-Spikes.
    \item \textbf{Staging-Umgebung}: Niemals direkt in Produktion deployen. Eine Staging-Umgebung mit identischer Architektur fängt die meisten Fehler.
    \item \textbf{Canary-Deployments}: Neue Prompt-Versionen nur für 5\% der User freischalten, dann hochskalieren.
\end{enumerate}

% ============================================================
% ANTI-PATTERNS
% ============================================================
\section{Anti-Patterns -- was du vermeiden solltest}

\begin{itemize}

    \item \textbf{Kein Timeout gesetzt} Ohne Timeout blockiert dein Request ewig und frisst Resources. Ein LLM-Request kann bei Überlast Minuten dauern. Immer 10--30s Timeout.
    \item \textbf{API-Ausfall ignoriert} Wenn OpenAI down ist, sollte deine App eine Fallback-Nachricht zeigen statt einem 500er Error. Graceful Degradation ist kein Nice-to-have.
    \item \textbf{Caching vergessen} Ein Redis-Cache ist kein Micro-Optimization, sondern ein fundamentales Architektur-Element. Ohne Cache zahlst du für jede identische Query neu.
    \item \textbf{Monolith ohne Trennung} LLM-Logik, Business-Logik und Datenbank-Zugriff in einer Datei. Das skaliert nicht. Trenne in Service-Layer.
    \item \textbf{Hardcodierte Prompt-Templates} Prompt-Templates im Source-Code bedeuten ein Re-Deploy für jede Text-Änderung. Nutze Config-Dateien oder eine Datenbank.
    \item \textbf{Fehlende Kosten-Alarme} Ohne Budget-Alarme kann ein einziger Bug oder ein Traffic-Spike tausende Euro verbrauchen. Setze Schwellwerte und Alarmierung.
\end{itemize}

% ============================================================
% PERFORMANCE
% ============================================================
\section{Performance und Skalierung}

\subsection{Horizontal Scaling}

LLM-Services sind typischerweise stateless und horizontal skalierbar:

\begin{lstlisting}[language=bash, caption=Docker Compose für skalierbaren LLM-Service]
version: "3.8"
services:
  llm-service:
    build: .
    environment:
      - OPENAI_API_KEY=${OPENAI_API_KEY}
      - REDIS_URL=redis://redis:6379
    deploy:
      replicas: 3
      resources:
        limits:
          cpus: "1.0"
          memory: "512M"

  redis:
    image: redis:7-alpine
    ports:
      - "6379:6379"

  prometheus:
    image: prom/prometheus
    volumes:
      - ./prometheus.yml:/etc/prometheus/prometheus.yml
    ports:
      - "9090:9090"

  grafana:
    image: grafana/grafana
    ports:
      - "3000:3000"
\end{lstlisting}

\subsection{Kosten-Kontrolle}

\begin{table}[H]
\centering
\begin{tabularx}{\linewidth}{lX}
\toprule
\textbf{Technik} & \textbf{Ersparnis} \\
\midrule
Caching mit Redis & 30--50\% \\
Modell-Tiering (Mini statt Full) & 80--90\% \\
Prompt-Optimierung (kürzere Prompts) & 20--40\% \\
Prompt-Caching (Anthropic/OpenAI Feature) & 50--90\% \\
Batching von Requests & 20--30\% \\
Context-Window-Management & 10--20\% \\
\bottomrule
\end{tabularx}
\caption{Kostenspar-Techniken für LLM-Produktionssysteme}
\end{table}

% ============================================================%
% CODE -- Kubernetes-Deployment
% ============================================================
\section{Codebeispiele -- Deployment-Optimierung}

\subsection{Multi-Region-Failover}

Ein produktionsreifes System überlebt einen Region- oder Provider-Ausfall:

\begin{lstlisting}[language=Python, caption={Multi-Region- und Multi-Provider-Failover}, label=lst:multiregion]
import random
from abc import ABC, abstractmethod

class LLMProvider(ABC):
    """Abstrakte Basis fur alle LLM-Provider."""

    @abstractmethod
    def complete(self, prompt: str) -> str:
        ...

class OpenAIProvider(LLMProvider):
    def complete(self, prompt: str) -> str:
        # Implementierung mit openai.Client()
        ...

class AnthropicProvider(LLMProvider):
    def complete(self, prompt: str) -> str:
        # Implementierung mit anthropic.Client()
        ...

class LocalProvider(LLMProvider):
    def complete(self, prompt: str) -> str:
        # Ollama / vLLM lokal
        ...

class FailoverRouter:
    """Router mit Multi-Provider-Failover."""

    def __init__(self, providers: list[LLMProvider]):
        self.providers = providers
        self.failed: set[int] = set()

    def complete(self, prompt: str,
                 max_retries: int = 3) -> str:
        for attempt in range(max_retries):
            for i, provider in enumerate(self.providers):
                if i in self.failed:
                    continue
                try:
                    return provider.complete(prompt)
                except Exception:
                    self.failed.add(i)
                    continue
        raise RuntimeError("Alle Provider ausgefallen")

# Nutzung
router = FailoverRouter([
    OpenAIProvider(),      # Primar: OpenAI
    AnthropicProvider(),   # Fallback 1: Anthropic
    LocalProvider(),       # Fallback 2: Lokal
])
\end{lstlisting}

\subsection{Health-Check-Endpoint mit Readiness-Probe}

\begin{lstlisting}[language=Python, caption={Health-Check mit API-Readiness}, label=lst:healthcheck]
from fastapi import FastAPI
import httpx

app = FastAPI()

@app.get("/health")
async def health():
    """Standard-Healthcheck."""
    return {"status": "ok"}

@app.get("/ready")
async def readiness():
    """Readiness-Check: Testet, ob die LLM-API erreichbar ist."""
    try:
        async with httpx.AsyncClient(timeout=5) as client:
            resp = await client.post(
                "https://api.openai.com/v1/chat/completions",
                headers={"Authorization": "Bearer sk-..."},
                json={
                    "model": "gpt-4o-mini",
                    "messages": [{"role": "user",
                                  "content": "ping"}],
                    "max_tokens": 1,
                },
            )
            if resp.status_code == 200:
                return {"status": "ready", "llm": "reachable"}
    except Exception:
        pass
    return {"status": "not_ready", "llm": "unreachable"}
\end{lstlisting}

\subsection{Kubernetes-Deployment-Manifest}

\begin{lstlisting}[language=Python, caption={Kubernetes-Deployment für LLM-Service (YAML)}]
# apiVersion: apps/v1
# kind: Deployment
# metadata:
#   name: llm-service
# spec:
#   replicas: 3
#   strategy:
#     rollingUpdate:
#       maxSurge: 1
#       maxUnavailable: 0
#   template:
#     spec:
#       containers:
#         - name: llm-service
#           image: llm-service:latest
#           env:
#             - name: OPENAI_API_KEY
#               valueFrom:
#                 secretKeyRef:
#                   name: llm-secrets
#                   key: openai-key
#           resources:
#             requests:
#               memory: "512Mi"
#               cpu: "500m"
#             limits:
#               memory: "1Gi"
#               cpu: "1000m"
#           livenessProbe:
#             httpGet: {path: /health, port: 8000}
#           readinessProbe:
#             httpGet: {path: /ready, port: 8000}
\end{lstlisting}

\praxishinweis{Kubernetes-Proben unterscheiden: livenessProbe = "Pod neu starten", readinessProbe = "Traffic stoppen". Ein LLM-Service mit Timeout sollte readiness-failen (Traffic umleiten), nicht liveness-failen (was einen Restart auslöst der Minuten dauert).}

% ============================================================
% SICHERHEIT
% ============================================================
\section{Sicherheit im Deployment}

\begin{enumerate}[leftmargin=*]
    \item \textbf{API-Keys sicher verwahren}: Nutze Secrets-Manager (AWS Secrets Manager, HashiCorp Vault, GitHub Secrets). Niemals Keys im Code oder in Docker-Images.
    \item \textbf{Network Isolation}: Der LLM-Service sollte nur über das API-Gateway erreichbar sein. Kein direkter Zugriff aus dem Internet.
    \item \textbf{Input Validation vor dem LLM-Call}: Validieren, dass Inputs nicht zu lang sind (max\_tokens), keine Injection-Versuche enthalten und kein PII unmaskiert übergeben wird.
    \item \textbf{Rate Limiting}: Schütze dich vor versehentlicher Kosten-Explosion und vor absichtlichem Abuse. Limits pro User/IP/API-Key.
    \item \textbf{Audit-Logging}: Jeder LLM-Call wird protokolliert: Wer hat was gefragt, welche Antwort kam, was hat es gekostet. Unverzichtbar für Compliance.
    \item \textbf{Dependency-Scanning}: Scan deine Dependencies regelmässig auf bekannte Sicherheitslücken (pip audit, Snyk, Dependabot).
\end{enumerate}

% ============================================================
% ZUSAMMENFASSUNG
% ============================================================
\begin{zusammenfassungEnv}
\begin{itemize}[leftmargin=*]
    \item FastAPI + Docker ist der Standard-Stack für LLM-Microservices
    \item Streaming und Caching sind die effektivsten Latenz-Kosten-Sparer
    \item Retry mit exponentiellem Backoff und Fallback-Modelle sind Pflicht
    \item Circuit-Breaker verhindert unnötige API-Calls bei Ausfällen
    \item Prometheus + Grafana + strukturiertes Logging = Observability
    \item Multi-Provider-Strategie verhindert Vendor-Lock-in
\end{itemize}
\end{zusammenfassungEnv}




\par\mbox{}\par
\begin{merkeEnv}
LLM-Deployment ist nicht wie klassisches Web-Deployment. External Dependencies, probabilistisches Verhalten und Kosten pro Request erfordern neue Patterns: Circuit-Breaker, Caching, Multi-Provider-Fallback und kontinuierliches Cost-Monitoring.
\end{merkeEnv}

% ============================================================
% PRAXISPROJEKT
% ============================================================
\section{Praxisprojekt -- Produktionsreifer LLM-Service mit Dashboard}

\textbf{Ziel:} Deploye einen LLM-Chat-Service mit FastAPI, Docker, Redis-Caching und Prometheus-Monitoring.

\textbf{Aufgaben:}
\begin{enumerate}[leftmargin=*]
    \item Erstelle einen FastAPI-Endpoint mit Rate Limiting (slowapi)
    \item Implementiere Redis-Caching mit TTL (30 Minuten)
    \item Füge Prometheus-Metriken für Requests, Latenz, Kosten hinzu
    \item Konfiguriere Docker Compose mit 3 Replicas des Services
    \item Richte Grafana-Dashboard ein: Latenz, Cost/Query, Error Rate
    \item Baue einen Circuit-Breaker für API-Ausfälle
\end{enumerate}

\textbf{Erweiterung:} Füge einen Multi-Provider-Fallback hinzu (OpenAI -> Anthropic -> lokales Modell).
Füge Readiness-Probes und Kubernetes-Manifeste für das Deployment hinzu.

% ============================================================
% RESSOURCEN
% ============================================================
\section{Ressourcen}

\begin{itemize}[leftmargin=*]
    \item \textbf{FastAPI Documentation}: \href{https://fastapi.tiangolo.com}{fastapi.tiangolo.com}
    \item \textbf{Prometheus \& Grafana}: prometheus.io, grafana.com -- Metriken und Dashboards
    \item \textbf{Tenacity}: \href{https://github.com/jd/tenacity}{github.com/jd/tenacity} -- Retry-Library für Python
    \item \textbf{Docker Compose}: \href{https://docs.docker.com/compose/}{docs.docker.com/compose}
    \item \textbf{OpenAI Production Best Practices}: \href{https://platform.openai.com/docs/guides/production-best-practices}{platform.openai.com}
    \item \textbf{LLM Deployment Patterns (Anthropic)}: \href{https://docs.anthropic.com/en/docs/quickstart}{docs.anthropic.com}
\end{itemize}

\textbf{Nächster Schritt:} Jetzt ist dein System in Produktion. Im nächsten Kapitel geht es um MLOps und Observability -- wie du es am Laufen hältst.
\endinput

\clearpage
